\documentclass[12pt]{article}
\usepackage{fullpage}
\usepackage{amsmath,amssymb,amsthm}

\newcommand{\cL}{\mathcal L}
\newcommand{\Dr}{\operatorname{Dr}}
\newcommand{\Trop}{\operatorname{Trop}}
\newcommand{\rk}{\operatorname{rk}}
\newcommand{\cl}{\operatorname{cl}}
\newcommand{\op}{\operatorname{op}}

\newtheorem{theorem}{Theorem}
\newtheorem{lemma}{Lemma}
\newtheorem{proposition}{Proposition}

\begin{document}

\begin{center}
{\bf Orthogonality on the relative Dressian of a self-opposite matroid}
\end{center}

\section*{Conventions}

Let $r=\rk M$ and let $\mu_M$ be the trivial rank $r$ valuated
matroid supported on the bases of $M$.  I use the convention of
Brandt--Eur--Zhang: a rank $k$ point of $\Dr(M)$ is a rank $k$
valuated matroid $\nu$ which is a valuated quotient of $\mu_M$, and
\[
        \Trop(\theta)\subseteq \Trop(\nu)
        \quad\Longleftrightarrow\quad
        \theta \text{ is a valuated quotient of }\nu .
\]
Write $p\preceq q$ for ``$p$ is a valuated quotient of $q$''.  Thus
$\preceq$ is the order opposite to inclusion of tropical linear spaces.
A projective tropical linear space determines its valuated matroid up
to adding a common scalar to all Plucker coordinates.

I will use two standard facts about valuated quotients.  First, for
consecutive ranks $d,d+1$, the quotient condition $p\preceq q$ is
equivalent, once $p$ and $q$ are valuated matroids, to the three-term
tropical incidence relations
\[
 \min\{p(Sx)+q(Syz),\,p(Sy)+q(Sxz),\,p(Sz)+q(Sxy)\}
 \text{ is attained at least twice}                       \tag{1}
\]
for all $|S|=d-1$ and distinct $x,y,z\notin S$.  In complete flags one
can use (1) without separately listing Plucker relations: a sequence
$g_0,\ldots,g_r$ with the trivial rank $0$ and rank $r$ endpoints is a
valuated complete flag if and only if all adjacent relations (1) hold.
This is the M-convex, or valuated generalized matroid, form of the
Dress--Wenzel--Murota cryptomorphisms.

Secondly, every valuated quotient $p\preceq q$ factors into elementary
quotients
\[
        p=h_a\preceq h_{a+1}\preceq\cdots\preceq h_b=q,
        \qquad \rk h_i=i.                                  \tag{2}
\]
This is the valuated Higgs-factorization theorem; in the case
$q=\mu_M$ the upper Higgs lifts are explicitly
\[
 h_j(B)=\min\{p(A): A\subseteq B,\ |A|=\rk p\}
\]
on independent $j$-subsets $B$ of $M$, with value $\infty$ otherwise.
These statements are the valuated versions of elementary quotient
factorization and follow from the M-convex/independent-set
cryptomorphisms of Dress--Wenzel and Murota; they are also equivalent
to the flag Dressian description of Brandt--Eur--Zhang.

Loops and parallel elements cause no additional difficulty.  The flat
lattice of $M$ is canonically the flat lattice of its simplification;
loops are loops in every quotient, and parallel elements have identical
coordinates in every quotient by the flat-constancy lemma below.  Hence
$\Dr(M)$ is the pullback of the relative Dressian of the simplification
by repeating parallel coordinates and adding loop coordinates.  We
therefore argue for the simple loopless matroid, and the formula
obtained is invariant under this pullback.

\section*{1. Consequences of the flat-lattice involution}

\begin{lemma}\label{lem:modular}
If the flat lattice $L=\cL(M)$ admits an order-reversing involution
$x\mapsto x^\perp$, then $L$ is modular and
$\rk(x^\perp)=r-\rk(x)$.
\end{lemma}

\begin{proof}
The lattice of flats is a finite ranked geometric lattice, hence upper
semimodular.  An anti-automorphism reverses maximal chains, so it sends
rank $i$ elements to rank $r-i$ elements.  Applying the
anti-automorphism to the semimodular inequality
\[
 \rk x+\rk y\ge \rk(x\wedge y)+\rk(x\vee y)
\]
gives the opposite inequality.  Hence equality holds for every pair
$x,y$, which is modularity.
\end{proof}

Thus $(x\vee y)^\perp=x^\perp\wedge y^\perp$ and
$(x\wedge y)^\perp=x^\perp\vee y^\perp$.

\begin{lemma}[flat constancy]\label{lem:flatconst}
Let $\nu\preceq \mu_M$ have rank $k$.  If $B,B'$ are independent
$k$-subsets of $M$ with $\cl(B)=\cl(B')$, then $\nu(B)=\nu(B')$.
Moreover $\nu(A)=\infty$ whenever $A$ is dependent in $M$.
\end{lemma}

\begin{proof}
The last assertion is the usual support statement for matroid
quotients: the underlying matroid of $\nu$ is a quotient of $M$, so
every basis of it is independent in $M$.

For $0<k<r$, the basis graph of the restriction $M|\cl(B)$ is
connected, so it suffices to compare adjacent bases
$B=A\cup\{b\}$ and $B'=A\cup\{b'\}$ with $|A|=k-1$ and
$\cl(Ab)=\cl(Ab')$.  Choose $C$ of size $r-k$ such that
$A\cup\{b\}\cup C$ is a basis of $M$; then
$A\cup\{b'\}\cup C$ is also a basis.  Put
$T=A\cup\{b,b'\}\cup C$.  In the incidence relation between
$\nu$ and $\mu_M$ with the set $A$ and the set $T$, the only finite
terms are the two terms $\nu(Ab)$ and $\nu(Ab')$: if $i\in C$, then
$T-i$ contains the dependent set $A\cup\{b,b'\}$.  The tropical
minimum condition forces $\nu(Ab)=\nu(Ab')$.  The cases $k=0$ and
$k=r$ are immediate: rank $0$ has one basis, while a rank $r$ quotient
of $\mu_M$ represents the same top tropical linear space and hence is
a scalar translate of $\mu_M$.
\end{proof}

Hence a quotient $\nu$ of rank $k$ has flat coordinates
\[
        \bar\nu(X)=\nu(B),\qquad \rk X=k,\quad \cl(B)=X.
\]

\section*{2. The local orthogonality check}

For a rank $k$ quotient $\nu\preceq\mu_M$ define a rank $r-k$
coordinate vector by
\[
 \nu^{\perp_M}(A)=
 \begin{cases}
   \bar\nu\bigl((\cl A)^\perp\bigr),
      & A\text{ independent in }M,\ |A|=r-k,\\
   \infty,&\text{otherwise.}
 \end{cases}                                                 \tag{3}
\]
It is visibly involutive on flat coordinates.  The key point is that
it carries elementary quotient steps to elementary quotient steps in
the opposite direction.

\begin{lemma}[elementary opposition]\label{lem:elementary}
Let $p\preceq q\preceq\mu_M$, where $\rk p=d$ and $\rk q=d+1$.
Then the two transformed coordinate functions
$q^{\perp_M}$ and $p^{\perp_M}$ satisfy every adjacent three-term
relation (1) of ranks $r-d-1$ and $r-d$.
\end{lemma}

\begin{proof}
Fix a set $A$ of size $r-d-2$ and
distinct elements $x,y,z\notin A$.  If $A$ is dependent, every term is
$\infty$, so assume $A$ is independent and put $U=\cl(A)$.  Let
$W=\cl(Axyz)$ and let $\delta=\rk W-\rk U$.

If $\delta\le 1$, each two-element set among $x,y,z$ is dependent
over $U$, so every $p^{\perp_M}$ term is $\infty$.

Assume $\delta=2$.  Let $G=W^\perp$ and $H=U^\perp$; then
$\rk G=d$ and $\rk H=d+2$.  Every finite term has the same
$p$-coordinate $\bar p(G)$.  The remaining coordinates are the
rank $d+1$ flats
\[
       L_x=(\cl(Ax))^\perp,\quad
       L_y=(\cl(Ay))^\perp,\quad
       L_z=(\cl(Az))^\perp
\]
in the rank-two interval $[G,H]$.  If two of these flats coincide, the
corresponding two finite terms are already equal.  If they are three
distinct points of the interval, choose a basis $S$ of $G$, choose
representatives $\ell_x,\ell_y,\ell_z$ for them in $H$, and extend
$S\cup\{\ell_x,\ell_y\}$ to a basis of $M$ by a set $C$ outside $H$.
The incidence relation for $q\preceq\mu_M$ with lower set $S$ and
upper set $S\cup\{\ell_x,\ell_y,\ell_z\}\cup C$ has precisely the
three finite terms
\[
       \bar q(L_x),\qquad \bar q(L_y),\qquad \bar q(L_z),
\]
and therefore their minimum is attained at least twice.  Adding the
common summand $\bar p(G)$ gives the desired relation.

It remains to consider $\delta=3$.  Now $G=W^\perp$ and $H=U^\perp$
have ranks $d-1$ and $d+2$.  In the rank-three interval $[G,H]$ set
\[
 P_x=(\cl(Ayz))^\perp,\quad P_y=(\cl(Axz))^\perp,\quad
 P_z=(\cl(Axy))^\perp,
\]
and
\[
 Q_x=(\cl(Ax))^\perp,\quad Q_y=(\cl(Ay))^\perp,\quad
 Q_z=(\cl(Az))^\perp .
\]
The flats $P_x,P_y,P_z$ are the three atoms over $G$ dual to the
independent frame $x,y,z$ over $U$, and
\[
        Q_x=P_y\vee P_z,\qquad
        Q_y=P_x\vee P_z,\qquad
        Q_z=P_x\vee P_y
\]
inside the interval $[G,H]$.  Choose a basis $S$ of $G$ and elements
$a_x\in P_x\setminus G$, $a_y\in P_y\setminus G$,
$a_z\in P_z\setminus G$.  The incidence relation (1) for the original
elementary quotient $p\preceq q$, applied to $S$ and
$a_x,a_y,a_z$, is exactly
\[
 \min\{\bar p(P_x)+\bar q(Q_x),\,
        \bar p(P_y)+\bar q(Q_y),\,
        \bar p(P_z)+\bar q(Q_z)\}
\]
attained at least twice.  This is the required relation for
$q^{\perp_M}\preceq p^{\perp_M}$.
\end{proof}

\section*{3. Construction of the involution}

\begin{theorem}
The map
\[
        \Phi(\Trop(\nu))=\Trop(\nu^{\perp_M})
\]
is a well-defined order-reversing involution of $\Dr(M)$.
\end{theorem}

\begin{proof}
Adding a scalar to $\nu$ adds the same scalar to $\nu^{\perp_M}$, so
the formula depends only on the projective tropical linear space.

Let $\nu\preceq\mu_M$.  By the valuated Higgs factorization, insert
$\nu$ into a complete flag
\[
  h_0\preceq h_1\preceq\cdots\preceq h_r=\mu_M,\qquad h_k=\nu .
\]
Applying Lemma \ref{lem:elementary} to each adjacent quotient shows
that the reversed sequence
\[
  h_r^{\perp_M},h_{r-1}^{\perp_M},\ldots,h_0^{\perp_M}
\]
satisfies all adjacent three-term relations.  Its endpoints are the
rank $0$ trivial valuated matroid and $\mu_M$, respectively; hence, by
the complete-flag criterion stated above, it is a valuated complete
flag.  Thus $h_k^{\perp_M}=\nu^{\perp_M}$ is a quotient of $\mu_M$, so
$\Phi$ maps $\Dr(M)$ to itself.  Since
$(X^\perp)^\perp=X$ for flats, (3) also gives
$(\nu^{\perp_M})^{\perp_M}=\nu$ up to scalar; hence $\Phi^2$ is the
identity.

Finally suppose $\Trop(\theta)\subseteq\Trop(\nu)$, equivalently
$\theta\preceq\nu$.  Factor the flag
$\theta\preceq\nu\preceq\mu_M$ into elementary steps and apply
Lemma \ref{lem:elementary} to each step.  The reversed transformed
complete sequence is again a valuated complete flag and contains
\[
        \nu^{\perp_M}\preceq\theta^{\perp_M}.
\]
By the quotient/inclusion convention this says
\[
        \Phi(\Trop(\nu))\subseteq\Phi(\Trop(\theta)),
\]
so $\Phi$ reverses inclusion.
\end{proof}

\section*{4. Agreement with the flat involution}

Let $F$ be a flat of rank $f$ and let
$q_F$ denote the rank $r-f$ valuated matroid of
$M/F\oplus U_{0,F}$.  A rank $r-f$ flat $Z$ has finite
$q_F$-coordinate exactly when $Z\vee F=\hat 1$.  For an independent
$f$-set $A$ with $Y=\cl(A)$,
\[
 \Phi(q_F)(A)<\infty
 \quad\Longleftrightarrow\quad
        Y^\perp\vee F=\hat 1 .
\]
Since $\rk(Y^\perp)=r-f$ and $\rk F=f$, modularity makes this
equivalent to $Y^\perp\wedge F=\hat 0$.  Applying $\perp$ gives
\[
        Y\vee F^\perp=\hat 1,
\]
which is precisely the condition that $A$ be a basis of
$M/F^\perp\oplus U_{0,F^\perp}$.  Dependent $A$ has value $\infty$ on
both sides.  Therefore
\[
 \Phi\bigl(\Trop(M/F\oplus U_{0,F})\bigr)
 =
 \Trop(M/F^\perp\oplus U_{0,F^\perp}),
\]
so $\Phi$ extends the given order-reversing involution on flats.

\begin{thebibliography}{9}

\bibitem{BEZ}
M.~Brandt, C.~Eur and L.~Zhang,
\emph{Tropical flag varieties},
Advances in Mathematics \textbf{384} (2021), 107695.

\bibitem{DW}
A.~Dress and W.~Wenzel,
\emph{Valuated matroids},
Advances in Mathematics \textbf{93} (1992), 214--250.

\bibitem{Murota}
K.~Murota,
\emph{Matroid valuation on independent sets},
Journal of Combinatorial Theory, Series B \textbf{69} (1997), 59--78.

\bibitem{Speyer}
D.~Speyer,
\emph{Tropical linear spaces},
SIAM Journal on Discrete Mathematics \textbf{22} (2008), 1527--1558.

\end{thebibliography}

\end{document}
